\subsection{Borel 子代数}

命题 7. 设 \(\mathfrak{b} = \mathfrak{h} + \mathfrak{g}^{\mathrm{P}}\) 是 \(\mathfrak{g}\) 的包含 \(\mathfrak{h}\) 的子代数。下列条件等价：

\begin{enumerate}
\def\labelenumi{(\roman{enumi})}
\item
  \(\mathfrak{b}\) 是 \(\mathfrak{g}\) 的极大可解子代数；
\item
  存在 R 的房 C 使 \(P = R_{+}(C)\) ；
\item
  \(\mathrm{P} \cap (-\mathrm{P}) = \varnothing\) 且 \(\mathrm{P} \cup (-\mathrm{P}) = \mathrm{R}\) 。
\item
  \(\Longrightarrow\) (ii)：若 \(\mathfrak{b}\) 可解，则 \(P \cap (-P) = \varnothing\) 。于是存在房 \(C\) （\(R\) 中的）使 \(P \subset R_{+}(C)\) （第六章 §1 第 7 节 命题 22）。这时 \(\mathfrak{h} + \mathfrak{g}^{R_{+}(C)}\) 是 \(\mathfrak{g}\) 的包含 \(\mathfrak{b}\) 的可解子代数，故若 \(\mathfrak{b}\) 极大则等于 \(\mathfrak{b}\) 。
\item
  \(\Longrightarrow\) (iii)：这是明显的。
\item
  \(\Longrightarrow\) (i)：设 \(P \cap (-P) = \varnothing\) 且 \(P \cup (-P) = R\) 。则 b 可解。设 \(b'\) 是 g 的包含 b 的可解子代数。存在 R 的子集 Q 使 \(b' = h + g^{Q}\) 。则 \(Q \cap (-Q) = \varnothing\) 且 \(Q \supset P\) ，故 Q = P 且 \(b' = b\) 。
\end{enumerate}

定义 1. g 的包含 h 且满足命题 7 中等价条件的子代数称为 (g, h) 的 Borel 子代数。

子代数 \(\mathfrak{b}\) （可分裂代数 \(\mathfrak{g}\) 的）称为 \(\mathfrak{g}\) 的 Borel 子代数，如果存在分裂嘉当子代数 \(\mathfrak{h}'\) （\(\mathfrak{g}\) 的）使 \(\mathfrak{b}\) 是 \((\mathfrak{g},\mathfrak{h}')\) 的 Borel 子代数。

设 \((\mathfrak{g},\mathfrak{h})\) 是分裂约化李代数。设 \(\mathfrak{g}=\mathfrak{c}\times\mathfrak{s}\) ，其中 c 交换而 s 半单。g 的形如 \(\mathfrak{c}\times\mathfrak{b}\) 的子代数（其中 b 是 \((\mathfrak{s},\mathfrak{h}\cap\mathfrak{s})\) 的 Borel 子代数）称为 \((\mathfrak{g},\mathfrak{h})\) 的 Borel 子代数。

采用命题 7 的记号，我们也说 b 是由 h 与 C（或由 h 与对应于 C 的 R 的基）所定义的 g 的 Borel 子代数。

注. 把 \(\mathrm{R}_{+}(\mathrm{C})\) 对应于 R 的每个房 C 的映射是单射（第六章 §1 第 7 节 命题 20 之推论 1）。因此，\(\mathrm{C} \mapsto \mathfrak{h} + \mathfrak{g}^{\mathrm{R}_{+}(\mathrm{C})}\) 是从 R 的房之集到 \((\mathfrak{g}, \mathfrak{h})\) 的 Borel 子代数之集上的双射。于是，\((\mathfrak{g}, \mathfrak{h})\) 的 Borel 子代数的个数等于 R 的 Weyl 群的阶（第六章 §1 第 5 节 定理 2）。

命题 8. 设 \(\mathfrak{b}\) 是 \(\mathfrak{g}\) 的子代数，\(k'\) 是 \(k\) 的扩张。则 \(\mathfrak{b} \otimes_k k'\) 是 \((\mathfrak{g} \otimes_k k', \mathfrak{h} \otimes_k k')\) 的 Borel 子代数当且仅当 \(\mathfrak{b}\) 是 \((\mathfrak{g}, \mathfrak{h})\) 的 Borel 子代数。

这由命题 7 的条件 (iii) 立即得出。

命题 9. 设 \(\mathfrak{b}\) 是由 R 的房 C 所定义的 \((\mathfrak{g},\mathfrak{h})\) 的 Borel 子代数。设 \(\mathfrak{n} = \mathfrak{g}^{\mathrm{R}_{+}(\mathrm{C})} = \sum_{\alpha \in \mathrm{R},\alpha >0}\mathfrak{g}^{\alpha}\) 。设 \(l = \dim \mathfrak{h}\) 。

\begin{enumerate}
\def\labelenumi{(\roman{enumi})}
\item
  若 \(h \in \mathfrak{h}\) 且 \(x \in \mathfrak{n}\) ，则 \(\mathrm{ad}_{\mathfrak{g}}(h + x)\) 的特征多项式是 \(\mathrm{T}^l \prod_{\alpha \in \mathbb{R}} (\mathrm{T} - \alpha(h))\) 。
\item
  \(\mathfrak{b}\) 的最大幂零理想等于 \(\mathfrak{n}\) ，也等于 \([\mathfrak{b},\mathfrak{b}]\) 。它也是 \(\mathfrak{b}\) 中在 \(\mathfrak{g}\) 内幂零的元素组成的集合。
\item
  设 \(\mathrm{B}\) 是 \(\mathbb{R}\) 的、对应于 \(\mathbb{C}\) 的基。对一切 \(\alpha \in \mathbb{B}\) ，设 \(X_{\alpha}\) 是 \(\mathfrak{g}^{\alpha}\) 的非零元素。则 \((X_{\alpha})_{\alpha \in \mathbb{B}}\) 生成李代数 \(\mathfrak{n}\) 。我们有 \([\mathfrak{n}, \mathfrak{n}] = \sum_{\alpha \in \mathbb{R}, \alpha > 0, \alpha \notin \mathbb{B}} \mathfrak{g}^{\alpha}\) 。
\end{enumerate}

在 \(\mathfrak{h}_{\mathbf{Q}}^{*}\) 上存在一个全序，它与它的向量空间结构相容，且使 \(\mathrm{R}_{+}(\mathrm{C})\) 的元素都 \(>0\) （第六章 §1 第 7 节）。设 \(h, x\) 如 (i)，并设 \(y \in \mathfrak{g}^{\alpha}\) 。则 \([h + x, y] = \alpha(h)y + z\) ，其中 \(z \in \sum_{\beta > \alpha} \mathfrak{g}^{\beta}\) 。于是，关于 \(\mathfrak{g}\) 的一个适当的基，\(\operatorname{ad}_{\mathfrak{g}}(h + x)\) 的矩阵具有下列性质：

\begin{enumerate}
\def\labelenumi{\arabic{enumi})}
\item
  它是下三角的；
\item
  矩阵的对角元是数 0（l 次）以及诸 \(\alpha(h)\) ，其中 \(\alpha \in \mathbb{R}\) 。
\end{enumerate}

这就证明了 (i)。它还表明 \(\operatorname{ad}_{\mathfrak{b}}(h+x)\) 的特征多项式是 \(\mathrm{T}^{l}\prod_{\alpha\in\mathrm{R}_{+}(\mathrm{C})}(\mathrm{T}-\alpha(h))\) 。由前面的论证，b 中在 g 内幂零的元素之集以及 b 的最大幂零理想都等于 n。~由命题 5 (i) 有 \(n=[b,b]\) 。最后，断言 (iii) 由 §2 命题 4 (ii) 以及第六章 §1 第 6 节 命题 19 得出。

推论. 设 \(\mathfrak{b}\) 是 \(\mathfrak{g}\) 的 Borel 子代数。

\begin{enumerate}
\def\labelenumi{(\roman{enumi})}
\item
  b 的每个嘉当子代数都是 g 的分裂嘉当子代数。
\item
  若 \(\mathfrak{h}_1, \mathfrak{h}_2\) 是 \(\mathfrak{b}\) 的嘉当子代数，则存在 \(x \in [\mathfrak{b}, \mathfrak{b}]\) 使 \(e^{\mathrm{ad}_{\mathfrak{g}} x} \mathfrak{h}_1 = \mathfrak{h}_2\) 。
\end{enumerate}

断言 (i) 由命题 9 (i) 与第七章 §3 第 3 节 命题 3 得出。断言 (ii) 由命题 9 (ii) 与第七章 §3 第 4 节 定理 3 得出。

命题 10. 设 \(\mathfrak{b},\mathfrak{b}'\) 是 \(\mathfrak{g}\) 的 Borel 子代数。则存在 \(\mathfrak{g}\) 的分裂嘉当子代数含于 \(\mathfrak{b} \cap \mathfrak{b}'\) 。

设 \(\mathfrak{h}\) 是 \(\mathfrak{b}\) 的嘉当子代数，\(\mathfrak{n} = [\mathfrak{b},\mathfrak{b}],\mathfrak{n}' = [\mathfrak{b}',\mathfrak{b}']\) ，\(\mathfrak{p} = \mathfrak{b}\cap \mathfrak{b}'\) ，并设 \(\mathfrak{s}\) 是 \(\mathfrak{g}\) 的、与 \(\mathfrak{b} + \mathfrak{b}'\) 互补的一个向量子空间。用 \(\mathfrak{s}^{\perp},\mathfrak{b}^{\perp},\mathfrak{b}'^{\perp}\) 表示 \(\mathfrak{s},\mathfrak{b},\mathfrak{b}^{\prime}\) 关于 \(\mathfrak{g}\) 的 Killing 型的正交补。置 \(l = \dim \mathfrak{h}, n = \dim \mathfrak{n}, p = \dim \mathfrak{p}\) 。则 \(\dim \mathfrak{b} = \dim \mathfrak{b}^{\prime} = l + n\) ，

\[
\dim \mathfrak {s} ^ {\perp} = \dim (\mathfrak {b} + \mathfrak {b} ^ {\prime}) = 2 (l + n) - p,
\]

于是

\[
\begin{array}{c} \dim (\mathfrak {s} ^ {\perp} \cap \mathfrak {p}) \geq \dim \mathfrak {s} ^ {\perp} + \dim \mathfrak {p} - \dim \mathfrak {g} \\ = 2 (l + n) - p + p - (l + 2 n) = l. \end{array}\tag{3}
\]

由 §2 第 2 节 命题 1，\(\mathfrak{n} \subset \mathfrak{b}^{\perp}\) ，\(\mathfrak{n}' \subset \mathfrak{b}'^{\perp}\) 。\(\mathfrak{p} \cap \mathfrak{n}\) 的元素在 \(\mathfrak{g}\) 内是幂零的（命题 9 (ii)），并且属于 \(\mathfrak{b}'\) ，从而属于 \(\mathfrak{n}'\) （命题 9 (ii)）。因此，\(\mathfrak{p} \cap \mathfrak{n} \subset \mathfrak{n} \cap \mathfrak{n}' \subset \mathfrak{b}^{\perp} \cap \mathfrak{b}'^{\perp}\) ，故 \(\mathfrak{s}^{\perp} \cap \mathfrak{p} \cap \mathfrak{n} = 0\) 。鉴于 (3)，我们看到 \(\mathfrak{s}^{\perp} \cap \mathfrak{p}\) 是 \(\mathfrak{n}\) 在 \(\mathfrak{b}\) 中的补。设 \(z\) 是 \(\mathfrak{h}\) 中在 \(\mathfrak{g}\) 内正则的一个元素；存在 \(y \in \mathfrak{n}\) 使 \(y + z \in \mathfrak{s}^{\perp} \cap \mathfrak{p}\) ；由命题 9 (i)，\(\operatorname{ad}_{\mathfrak{g}}(y + z)\) 与 \(\operatorname{ad}_{\mathfrak{g}}z\) 有相同的特征多项式，故 \(x = y + z\) 在 \(\mathfrak{g}\) 中正则，从而更有在 \(\mathfrak{b}\) 中与 \(\mathfrak{b}'\) 中正则（第七章 §2 第 2 节 命题 9）。由于 \(\mathfrak{g}, \mathfrak{b}, \mathfrak{b}'\) 有相同的秩，\(\mathfrak{b}^{0}(x) = \mathfrak{g}^{0}(x) = \mathfrak{b}'^{0}(x)\) 同时是 \(\mathfrak{b}\) 的、\(\mathfrak{g}\) 的与 \(\mathfrak{b}'\) 的嘉当子代数（第七章 §3 第 3 节 定理 2）。最后，\(\mathfrak{g}\) 的这个嘉当子代数由命题 9 的推论是分裂的。

推论. 群 \(\operatorname{Aut}_{e}(\mathfrak{g})\) 传递地作用于偶 \((\mathfrak{t},\mathfrak{b})\) 组成的集合上，其中 t 是 g 的分裂嘉当子代数，b 是 \((\mathfrak{g},\mathfrak{t})\) 的 Borel 子代数。

设 \((\mathfrak{t}_1, \mathfrak{b}_1)\) 与 \((\mathfrak{t}_2, \mathfrak{b}_2)\) 是两个这样的偶。存在分裂嘉当子代数 \(\mathfrak{t}\) （\(\mathfrak{g}\) 的）含于 \(\mathfrak{b}_1 \cap \mathfrak{b}_2\) （命题 10）。由命题 9 的推论，我们归结到 \(\mathfrak{t}_1 = \mathfrak{t}_2 = \mathfrak{t}\) 的情形。设 \(S\) 是 \((\mathfrak{g}, \mathfrak{t})\) 的根系。存在基 \(B_1, B_2\) （\(S\) 的）使 \(\mathfrak{b}_i\) 对应于 \(B_i\) （ \(i = 1, 2\) ），且存在 \(s \in W(S)\) 把 \(B_1\) 变为 \(B_2\) 。最后，存在 \(a \in \operatorname{Aut}_e(\mathfrak{g})\) 使 \(a|t = s\) （ \(\S 2\) 第 2 节 定理 2 的推论）。则 \(a(t) = t\) 且 \(a(\mathfrak{b}_1) = \mathfrak{b}_2\) 。

